% ============================================================
\subsubsection{Planificación de requisitos}
% ============================================================

\paragraph{Requisitos funcionales y no funcionales}

Los requisitos se obtuvieron a partir del levantamiento de información con el
actor principal del sistema, el nutricionista.
\pendiente{indica cómo se levantaron: reunión/entrevista con quién, fecha y si
fueron validados por el director o un nutricionista}.
La Tabla~\ref{tab:rf} presenta los doce requisitos funcionales y la
Tabla~\ref{tab:rnf} los siete requisitos no funcionales.

{\small
\begin{longtable}{|p{1.3cm}|p{12.2cm}|}
\caption{Requisitos funcionales del prototipo NutriAPP}\label{tab:rf}\\
\hline
\textbf{Código} & \textbf{Descripción} \\ \hline
\endfirsthead
\hline
\textbf{Código} & \textbf{Descripción} \\ \hline
\endhead
RF01 & El sistema debe permitir al nutricionista registrar un nuevo paciente con sus datos personales básicos (nombre, apellido, correo electrónico, teléfono, fecha de nacimiento y sexo). \\ \hline
RF02 & El sistema debe permitir al nutricionista registrar y actualizar el historial clínico del paciente, incluyendo antecedentes médicos, quirúrgicos y familiares, motivo de consulta, medicamentos, alergias e intolerancias alimentarias. \\ \hline
RF03 & El sistema debe permitir al nutricionista registrar medidas antropométricas del paciente: peso, talla, circunferencia de cintura y cadera, porcentaje de grasa corporal y masa muscular. \\ \hline
RF04 & El sistema debe calcular automáticamente el Índice de Masa Corporal (IMC) y la relación cintura-cadera, mostrando los resultados al nutricionista para su análisis clínico. \\ \hline
RF05 & El sistema debe permitir al nutricionista registrar datos biométricos del paciente: glucosa, HbA1c, colesterol total, triglicéridos, HDL, LDL, presión arterial y hemoglobina. \\ \hline
RF06 & El sistema debe permitir al nutricionista registrar los hábitos dietéticos del paciente, incluyendo frecuencia de comidas, consumo de agua, preferencias y restricciones alimentarias. \\ \hline
RF07 & El sistema debe calcular y proporcionar al nutricionista el Gasto Energético Total (GET) del paciente mediante la fórmula de Harris-Benedict, considerando sexo, peso, talla, edad y factor de actividad física seleccionado. \\ \hline
RF08 & El sistema debe calcular y proporcionar al nutricionista la distribución de macronutrientes diarios (proteínas, carbohidratos y grasas) en gramos y porcentajes, visualizando los resultados mediante un gráfico de dona interactivo. \\ \hline
RF09 & El sistema debe permitir al nutricionista interactuar con un asistente de inteligencia artificial mediante un chat individual por paciente, utilizando los datos clínicos, antropométricos y biométricos del paciente seleccionado como contexto. \\ \hline
RF10 & El sistema debe permitir al nutricionista consultar un chat global con el asistente de inteligencia artificial, el cual tiene acceso al contexto clínico de todos los pacientes registrados en el sistema. \\ \hline
RF11 & El sistema debe permitir al nutricionista cargar documentos PDF como base de conocimiento del asistente de inteligencia artificial, mediante un pipeline de generación aumentada por recuperación (RAG). \\ \hline
RF12 & El sistema debe permitir al nutricionista gestionar el catálogo de alimentos ecuatorianos almacenado en la base de datos, incluyendo nombre, grupo alimentario y composición nutricional por cada 100 gramos. \\ \hline
\end{longtable}
}

{\small
\begin{longtable}{|p{1.5cm}|p{2.6cm}|p{9.4cm}|}
\caption{Requisitos no funcionales del prototipo NutriAPP}\label{tab:rnf}\\
\hline
\textbf{Código} & \textbf{Categoría} & \textbf{Descripción} \\ \hline
\endfirsthead
\hline
\textbf{Código} & \textbf{Categoría} & \textbf{Descripción} \\ \hline
\endhead
RNF01 & Rendimiento & El sistema debe procesar y responder las consultas al asistente de inteligencia artificial en un tiempo máximo de 10 segundos bajo condiciones normales de conectividad. \\ \hline
RNF02 & Seguridad & El sistema debe garantizar la protección de las credenciales del nutricionista mediante cifrado de contraseñas con bcrypt y gestión de sesiones mediante tokens JWT en cabeceras HTTP. \\ \hline
RNF03 & Precisión & El margen de error en los cálculos automáticos de indicadores corporales (IMC, relación cintura-cadera) y gasto energético (Harris-Benedict) no debe exceder el 1\,\% respecto a los valores de referencia establecidos por las fórmulas oficiales. \\ \hline
RNF04 & Usabilidad & El sistema debe ofrecer una interfaz web clara e intuitiva, accesible desde navegadores modernos (Chrome, Firefox, Edge), sin necesidad de instalación adicional por parte del nutricionista. \\ \hline
RNF05 & Portabilidad & El sistema debe poder ejecutarse en un entorno de desarrollo local sin dependencia de infraestructura en la nube, utilizando únicamente la API de Groq para la inferencia del modelo de lenguaje durante la fase de prototipado. \\ \hline
RNF06 & Privacidad & Los datos clínicos procesados por el sistema corresponden a perfiles de prueba construidos con fines experimentales; en una implementación productiva, el sistema deberá cumplir con la Ley Orgánica de Protección de Datos Personales, garantizando el consentimiento explícito y el cifrado de la información del paciente. \\ \hline
RNF07 & Mantenibilidad & El código fuente del sistema debe estar versionado en GitHub con ramas \texttt{main} y \texttt{develop}, garantizando la trazabilidad del desarrollo y facilitando futuras actualizaciones. \\ \hline
\end{longtable}
}

\paragraph{Actores}

Se identificó un único actor humano, descrito en la Tabla~\ref{tab:actores}.
La API de Groq interviene como sistema externo de inferencia, no como actor
que inicie casos de uso.

\begin{table}[H]
\centering\small
\caption{Actores identificados}\label{tab:actores}
\begin{tabularx}{\linewidth}{|p{2.6cm}|X|}
\hline
\textbf{Actor} & \textbf{Descripción} \\ \hline
Nutricionista & Profesional de la salud que utiliza el sistema para registrar y gestionar los datos clínicos, antropométricos y biométricos de sus pacientes, consultar el asistente de inteligencia artificial y visualizar recomendaciones de planes nutricionales personalizados. \\ \hline
\end{tabularx}
\end{table}

\paragraph{Modelo conceptual}

A partir de los requisitos funcionales se elaboró el modelo conceptual de la
Figura~\ref{fig:conceptual}, que identifica las entidades principales del
dominio y sus relaciones. La Tabla~\ref{tab:glosario} define cada término.

\figura{modelo_conceptual.png}{Modelo conceptual de NutriAPP}{fig:conceptual}

\begin{table}[H]
\centering\small
\caption{Glosario de términos del modelo conceptual}\label{tab:glosario}
\begin{tabularx}{\linewidth}{|p{3.8cm}|X|}
\hline
\textbf{Término} & \textbf{Definición} \\ \hline
Nutricionista & Usuario autenticado que registra pacientes, valida planes y conversa con el asistente de IA. \\ \hline
Paciente & Persona atendida por el nutricionista; no posee credenciales de acceso al sistema. \\ \hline
HistorialClinico & Antecedentes médicos, quirúrgicos y familiares, motivo de consulta, medicación, alergias e intolerancias. \\ \hline
MedidasAntropometricas & Registro fechado de peso, talla, circunferencias, \% de grasa y masa muscular. \\ \hline
DatosBiometricos & Registro fechado de valores bioquímicos (glucosa, HbA1c, perfil lipídico, hemoglobina) y presión arterial. \\ \hline
HabitosDieteticos & Frecuencia de comidas, consumo de agua, preferencias y restricciones alimentarias. \\ \hline
SessionChat & Conversación del nutricionista con el asistente, individual por paciente o global. \\ \hline
MensajeChat & Mensaje del nutricionista o respuesta del asistente dentro de una sesión. \\ \hline
DocumentoConocimiento & Archivo PDF cargado por el nutricionista como contexto adicional del asistente. \\ \hline
PlanNutricional & Recomendación generada por el asistente y validada por el nutricionista. \\ \hline
Alimento & Elemento del catálogo con su grupo de equivalentes y composición nutricional. \\ \hline
\end{tabularx}
\end{table}
